%% ==========================================================================
%%  3  Our Results
%%
%%  Overview section.  It states the two conjectures --- which have no home
%%  elsewhere, since nothing proves them --- and points at where each result
%%  is proved.  Full statements and proofs live in the sections that produce
%%  them, so that nothing is numbered twice.
%% ==========================================================================
\section{Our Results}
\label{sec:results}

The target is Conjecture~\ref{conj:main}, stated in \Cref{sec:intro}: an
envy-free solution with $\subsidy \in \set{0,1}^n$ and hence total subsidy at
most $n-1$. Example~\ref{ex:lowerbound} supplies the matching lower bound, so it
would be tight. By Observation~\ref{obs:integrality} the conjecture is
equivalent to the purely graph-theoretic statement that some allocation $\alloc$
has no positive-weight cycle in $\envygraph{\alloc}$ and satisfies
$\pathw{\alloc}(i) \le 1$ for every agent $i$.

We do not settle Conjecture~\ref{conj:main}. What we contribute is a map of the
problem: two solved special cases, three structural tools, and --- the substance
of the paper --- two theorems showing that the two most natural proof strategies
both fail, one of them the strategy that succeeds on the goods side.

\paragraph{Solved cases.}
Conjecture~\ref{conj:main} holds, constructively and in polynomial time, for
\textbf{binary additive} costs (Theorem~\ref{thm:binadd}) and for
\textbf{identical} costs (Theorem~\ref{thm:identical}). The case $n = 2$ follows
from~\cite{BKNS22} through the complementation of Remark~\ref{rem:complement}.
The binary additive proof is the one to imitate: dump every chore on somebody
who finds it free, then balance whatever nobody wants. Its bound telescopes
along paths in the envy graph, which is the mechanism all three tools below try
to reproduce.

\paragraph{Structural tools} (\Cref{sec:structure}).
The arc-update table of Lemma~\ref{lem:arcupdate} isolates the entire difference
between the goods and chore problems in two lines. Proposition~\ref{prop:twotier}
characterises $\set{0,1}$ solutions as a two-tier structure, which is the natural
certificate format and suggests choosing the tiers before allocating.
Theorem~\ref{thm:sizeshift} exhibits the chore problem as \emph{exactly} the
goods problem of~\cite{BKNS22} plus a cardinality-balancing requirement, which
is the precise sense in which Conjecture~\ref{conj:main} is not a corollary
of~\cite{BKNS22}.

\paragraph{First negative result} (\Cref{sec:approach1}).
Item-by-item insertion --- allocate one chore at a time, never move an
already-allocated chore, and repair only by permuting whole bundles --- is the
template of~\cite{BKNS22}. Theorem~\ref{thm:obstruction} exhibits a three-agent
instance on which that template provably cannot be completed: a legal partial
state with $\subsidy \in \set{0,1}^3$ and an unallocated chore such that
\emph{every} agent who receives it forces a subsidy of $2$, while the full
instance admits a zero-subsidy allocation. Any proof of
Conjecture~\ref{conj:main} must therefore be allowed to re-open bundles that
have already been allocated. We state this bluntly because the template is
seductive and patching the insertion invariant is a dead end, not an unsolved
exercise.

\paragraph{Second negative result} (\Cref{sec:approach2}).
The natural repair is to stop building the allocation incrementally and instead
select a good allocation globally, the obvious candidate being one that
minimises total cost. This suggests the following strengthening.

\begin{conjecture}
\label{conj:msw}
The allocation in Conjecture~\ref{conj:main} may in addition be taken
\emph{utilitarian-optimal}, i.e.\ minimising $\sum_{i \in \agents}
\cost_i(\bundle{i})$ over all allocations.
\end{conjecture}

Proposition~\ref{prop:msw-false} refutes Conjecture~\ref{conj:msw}: there is a
three-agent, four-chore instance whose \emph{unique} utilitarian-optimal
allocation requires a subsidy of $2$, although another allocation --- of
strictly higher total cost --- requires no subsidy at all. Utilitarian
optimality is therefore not merely an unhelpful restriction but an actively
wrong one, and the programme of searching for a tie-breaking rule inside the
utilitarian-optimal set is unsound as posed.

\paragraph{A third approach, still open} (\Cref{sec:approach3}).
Both failures above are failures of \emph{timing}: insertion commits a chore to
an owner and cannot take it back, and utilitarian optimality commits to the
wrong allocation outright. \Cref{sec:approach3} changes coordinates so that
ownership is decided only by the last move that touches a chore.
Theorem~\ref{thm:replica} re-expresses the chore instance as a \emph{goods}
instance on an item set carrying $n-1$ copies of each chore, with identical envy
graphs, so the dichotomous class is closed under the duality and~\cite{BKNS22}
applies to the replica verbatim. Corollary~\ref{cor:coverage} then reduces
Conjecture~\ref{conj:main} to a single residual constraint --- \emph{coverage},
that no agent receive two copies of one type. Read dynamically this is the
\emph{peel} process, in which the orientation of~\cite{BKNS22} is restored
because the quantity handed out is relief rather than burden, and the witness of
Theorem~\ref{thm:obstruction} is handled with the invariant intact. The approach
is not known to work: Proposition~\ref{prop:deadends} exhibits legal states from
which no legal continuation exists, so what remains is to find a peel rule that
never enters one.

\paragraph{A third negative result} (\Cref{sec:approach4}).
Given Corollary~\ref{cor:coverage}, the natural move is to change the numbers so
that a coverage violation becomes unattractive, and thereby reuse~\cite{BKNS22}
as a black box. \Cref{sec:approach4} closes that at both levels. No dichotomous
reweighting of the replica instance separates coverage from non-coverage
(Theorems~\ref{thm:w4} and~\ref{thm:w4prime}, the latter exhaustive over all
$990^3$ candidate triples), and no inflation of item weights can steer the
algorithm of~\cite{BKNS22}, because its repair subroutine provably runs where
every candidate's marginal is $0$ and a duplicate moves no arc at all
(Theorem~\ref{thm:f2}, Corollary~\ref{cor:f3}). Both failures have one cause,
named in \Cref{sec:holderblind}: a rule fixed before the allocation is chosen
cannot act on a quantity the allocation determines. Coverage must be enforced by
the procedure that builds the allocation, not by the objective it is scored
against.

\paragraph{A working algorithm} (\Cref{sec:approach6}).
The size-shift transform of Theorem~\ref{thm:sizeshift} is an involution
(Lemma~\ref{lem:sizeshift-bijection}), so the chore problem restates as a pure
\emph{goods} question, Target G, with no chores or scheduling in it
(Proposition~\ref{prop:targetG-implies}). Running~\cite{BKNS22}'s own algorithm
on that goods instance fails for a structural reason --- item-by-item insertion
can lock two items into the same bundle permanently, verified by an exhaustive
backtracking search that finds a $0$-spread allocation \emph{unreachable} by any
legal execution of the algorithm (\Cref{sec:insertion-fails-goods}). Abandoning
insertion for a construct-and-repair procedure --- one pass of iterated
maximum-weight perfect matching, repaired by cardinality-preserving single-item
swaps, restarted on failure --- has \textbf{zero failures across $389{,}215$
test instances}, including every hard instance that defeated
Approaches~1 through~5.

Four properties of that algorithm are now \emph{theorems}: its outputs are
certified (Theorem~\ref{thm:algorithm-soundness}), it terminates in polynomial
time $O(n^3 m^3 \tau + n^4 m^3)$ (Theorem~\ref{thm:algorithm-complexity}), its
search space is genuinely the unordered balanced partitions
(Lemma~\ref{lem:assignment-invariance}), and --- the main theoretical result of
this line --- \textbf{it is complete for two agents}
(Theorem~\ref{thm:n2-complete}): every two-agent dichotomous instance admits a
balanced split of spread at most $1$, found constructively in $O(m\tau)$ time
with no restarts, by a discrete intermediate-value argument on the aggregate
$u = \tilde{v}_1 + \tilde{v}_2$ along swap paths. This strengthens even the
known $n = 2$ chore result, adding balancedness that the complementation route
does not provide. For $n \ge 3$, completeness remains the one open statement
(\Cref{sec:construct-repair-remains}).

\paragraph{Evidence} (\Cref{sec:experiments}).
Conjecture~\ref{conj:main} survives exhaustive verification for $n = m = 3$ over
all $9{,}880$ instances up to agent symmetry, and randomised search over some
$3.5 \times 10^4$ further instances. \Cref{sec:experiments} also explains why
that evidence is weaker than the counts suggest, and what would strengthen it.
